\documentclass[12pt]{article}
\usepackage{amsmath}
\usepackage{amsfonts}
\usepackage{float}
\usepackage{fancyhdr}
\usepackage{graphicx}
\usepackage{setspace}
\usepackage[colorlinks=true,linkcolor=blue, citecolor=red]{hyperref}
\usepackage{url}
\usepackage[top=.75in, left=.75in, right=.75in, bottom=1in]{geometry}

% For algorithm
\usepackage{algorithm}
\usepackage{algpseudocode}

% ============ CODE ============
\usepackage{listings}
\usepackage{xcolor}
\definecolor{codegreen}{rgb}{0,0.6,0}
\definecolor{codegray}{rgb}{0.5,0.5,0.5}
\definecolor{codepurple}{rgb}{0.58,0,0.82}
\definecolor{backcolour}{rgb}{0.95,0.95,0.92}

% Styling for the code.
\lstdefinestyle{mystyle}{
    backgroundcolor=\color{backcolour},   
    commentstyle=\color{codegreen},
    keywordstyle=\color{magenta},
    numberstyle=\tiny\color{codegray},
    stringstyle=\color{codepurple},
    basicstyle=\ttfamily\footnotesize,
    breakatwhitespace=false,         
    breaklines=true,                 
    captionpos=b,                    
    keepspaces=true,                 
    numbers=left,                    
    numbersep=5pt,                  
    showspaces=false,                
    showstringspaces=false,
    showtabs=false,                  
    tabsize=2
}
\lstset{style=mystyle}

% Table packages and Unicode / Font support
\usepackage{array}
\usepackage{longtable}
\usepackage{enumitem}
\usepackage{booktabs}
\usepackage[utf8]{inputenc}
\DeclareUnicodeCharacter{2264}{\ensuremath{\le}}
\DeclareUnicodeCharacter{2194}{\ensuremath{\leftrightarrow}}
\usepackage[T1]{fontenc}
\usepackage{lmodern}
\graphicspath{{../../../../templates/}{../../prototype/}{figures/}{./}}
\setlength{\emergencystretch}{3em}
\setlist[enumerate]{label=\arabic*.,leftmargin=*}
\setlist[itemize]{leftmargin=*}

% ============ HEADER AND FOOTER ============
\setlength{\headheight}{29.43912pt}

\pagestyle{fancy}
\fancyfoot{}
\fancyfoot[R]{Page \thepage}

\setlength{\parindent}{0pt}
\renewcommand{\baselinestretch}{1.5}

\newcommand{\coursename}{Introduction of Software Engineering}
\newcommand{\reportname}{Studitopia}
\newcommand{\reporttitle}{\href{https://github.com/RoundFloor/Studitopia.git}{Studitopia}}

\newcommand{\studentname}{%
Nguyen Thanh An (24125024)\\
Nguyen Chi Bao (24125025)\\
Tran Nguyen Khanh Binh (24125026)\\
Tran Canh Anh Tuan (24125048)\\
Nguyen Le Quoc Huy (24125100)
}
\newcommand{\teachername}{Ho Tuan Thanh, Nguyen Thi Minh Tuyen}

% Header
\lhead{\reporttitle}
\rhead{
University of Science - VNUHCM\\
\coursename
}

% Footer
\newcommand{\leftfooter}{}
\lfoot{\leftfooter}

% ============ DOCUMENT ============
\begin{document}

\pagenumbering{roman}
\begin{titlepage}
\newcommand{\HRule}{\rule{\linewidth}{0.5mm}}
\centering

\textsc{\LARGE VNU-HCM}\\[0.5cm]
\textsc{\Large University of Science}\\[0.5cm]
\textsc{\large Faculty of Information Technology}\\[0.5cm]
\textsc{Computer Science}\\[0.5cm]

\HRule \\[0.4cm]
{ 
\huge{\bfseries{\reporttitle}}\\[0.5cm]
\large{\bfseries{UC-33 --- Control Communication}}
}\\[0.4cm]
\HRule \\[0.5cm]

\textbf{\large Subject: \coursename}\\[0.5cm]

\begin{minipage}[t]{0.55\textwidth}
\begin{flushleft} \large
\emph{Students:}\\
\studentname
\end{flushleft}
\end{minipage}
~
\begin{minipage}[t]{0.4\textwidth}
\begin{flushright} \large
\emph{Supervising Teachers:} \\
\teachername
\end{flushright}
\end{minipage}\\[1cm]

{\large \today}\\[1cm]

\includegraphics[scale=.20]{img/hcmus-logo.png}\\[1cm] 

\vfill
\end{titlepage}	

\tableofcontents
\pagebreak

\pagenumbering{arabic}
\setcounter{page}{1}

\hypertarget{sec:uc-33-control-communication}{}
\addcontentsline{toc}{section}{UC-33 --- Control Communication}
\section*{UC-33 --- Control Communication}
\par\medskip

\textbf{Version:} \texttt{1.0}
\par\medskip

\textbf{Performed by:} [Tran Canh Anh Tuan] | \textbf{Edited by:} [Tran Canh Anh Tuan]
\par\medskip

\begingroup
\small
\renewcommand{\baselinestretch}{1.0}\selectfont
\renewcommand{\arraystretch}{1.15}
\setlength{\tabcolsep}{6pt}
\begin{longtable}{|>{\raggedright\arraybackslash}p{0.09\textwidth}|>{\raggedright\arraybackslash}p{0.84\textwidth}|}
\hline
\textbf{Field} & \textbf{Value} \\ \hline
\endhead
Use Case ID & UC-33 \\ \hline
Name & Control Communication \\ \hline
Primary Actor & User \\ \hline
Supporting Actors & None \\ \hline
Feature Group & Notifications and Safety \\ \hline
Priority & Not explicitly stated \\ \hline
Related User Stories & US-39 \\ \hline
Related Requirements & AC-39.1, AC-39.2, AC-39.3; FR-39.1, FR-39.2, FR-39.3, FR-39.4, FR-39.5; NFR-SEC-1, NFR-SEC-3, NFR-SEC-6, NFR-SEC-8, NFR-SEC-9, NFR-PRIV-2, NFR-PRIV-6, NFR-PRIV-7 \\ \hline
Source Documents & \texttt{proposals/userstory.md} (sections 11.2-11.3); \texttt{proposals/functional\_requirement.md} (US-39); \texttt{proposals/settings.md}; \texttt{proposals/friend.md}; \texttt{proposals/friend\_profile.md}; \texttt{proposals/group\_detail.md}; \texttt{proposals/moderation\_queue.md}; \texttt{proposals/non\_functional\_requirement.md}; \texttt{proposals/API.md}; \texttt{proposals/response\_code.md}; \texttt{proposals/feature\_flow.md}; \texttt{proposals/screen\_flow.md} \\ \hline
\end{longtable}
\endgroup
\par\medskip

\hypertarget{sec:part-a-use-case-model}{}
\addcontentsline{toc}{subsection}{Part A --- Use Case Model}
\subsection*{Part A --- Use Case Model}
\par\medskip

\hypertarget{sec:a1-focused-use-case-diagram}{}
\addcontentsline{toc}{subsubsection}{A1. Focused Use Case Diagram}
\subsubsection*{A1. Focused Use Case Diagram}
\par\medskip

\begin{center}
\includegraphics[width=0.85\textwidth,height=0.75\textheight,keepaspectratio]{"figures/UC-33-control-communication_1.png"}
\end{center}
\par\medskip

\textbf{Diagram notation.} Person shape = human actor; cloud = supporting system; blue = selected use case; solid line = association; labelled solid arrow = generalization; dashed arrow = include or extend.
\par\medskip

\hypertarget{sec:a2-actor-register}{}
\addcontentsline{toc}{subsubsection}{A2. Actor Register}
\subsubsection*{A2. Actor Register}
\par\medskip

\begingroup
\small
\renewcommand{\baselinestretch}{1.0}\selectfont
\renewcommand{\arraystretch}{1.15}
\setlength{\tabcolsep}{6pt}
\begin{longtable}{|>{\raggedright\arraybackslash}p{0.09\textwidth}|>{\raggedright\arraybackslash}p{0.09\textwidth}|>{\raggedright\arraybackslash}p{0.74\textwidth}|}
\hline
\textbf{Actor} & \textbf{Type} & \textbf{Description} \\ \hline
\endhead
User & Primary & A registered user who controls communication and privacy settings so that other users cannot interact in unpermitted ways (US-39). \\ \hline
\end{longtable}
\endgroup
\par\medskip

\hypertarget{sec:a3-relationships}{}
\addcontentsline{toc}{subsubsection}{A3. Relationships}
\subsubsection*{A3. Relationships}
\par\medskip

\hypertarget{sec:a3-1-associations}{}
\paragraph*{A3.1 Associations}
\par\medskip

\begingroup
\small
\renewcommand{\baselinestretch}{1.0}\selectfont
\renewcommand{\arraystretch}{1.15}
\setlength{\tabcolsep}{6pt}
\begin{longtable}{|>{\raggedright\arraybackslash}p{0.09\textwidth}|>{\raggedright\arraybackslash}p{0.15\textwidth}|>{\raggedright\arraybackslash}p{0.09\textwidth}|>{\raggedright\arraybackslash}p{0.58\textwidth}|}
\hline
\textbf{Actor} & \textbf{Use Case} & \textbf{Type} & \textbf{Description} \\ \hline
\endhead
User & UC-33 Control Communication & Primary & Opens Settings and configures permitted interactions, restrictions, muted members, and blocked members. \\ \hline
\end{longtable}
\endgroup
\par\medskip

\hypertarget{sec:a3-2-generalization-and-specialization}{}
\paragraph*{A3.2 Generalization and Specialization}
\par\medskip

\begingroup
\small
\renewcommand{\baselinestretch}{1.0}\selectfont
\renewcommand{\arraystretch}{1.15}
\setlength{\tabcolsep}{6pt}
\begin{longtable}{|>{\raggedright\arraybackslash}p{0.09\textwidth}|>{\raggedright\arraybackslash}p{0.09\textwidth}|>{\raggedright\arraybackslash}p{0.09\textwidth}|>{\raggedright\arraybackslash}p{0.65\textwidth}|}
\hline
\textbf{Child} & \textbf{Parent} & \textbf{Type} & \textbf{Inherited Behavior or Association} \\ \hline
\endhead
None & None & None & The proposals do not name an actor generalization or use-case specialization for US-39. \\ \hline
\end{longtable}
\endgroup
\par\medskip

\hypertarget{sec:a3-3-include}{}
\paragraph*{A3.3 Include}
\par\medskip

\begingroup
\small
\renewcommand{\baselinestretch}{1.0}\selectfont
\renewcommand{\arraystretch}{1.15}
\setlength{\tabcolsep}{6pt}
\begin{longtable}{|>{\raggedright\arraybackslash}p{0.13\textwidth}|>{\raggedright\arraybackslash}p{0.16\textwidth}|>{\raggedright\arraybackslash}p{0.63\textwidth}|}
\hline
\textbf{Base Use Case} & \textbf{Included Use Case} & \textbf{Meaning} \\ \hline
\endhead
None & None & The proposals do not name a distinct included use case for US-39. \\ \hline
\end{longtable}
\endgroup
\par\medskip

\hypertarget{sec:a3-4-extend}{}
\paragraph*{A3.4 Extend}
\par\medskip

\begingroup
\small
\renewcommand{\baselinestretch}{1.0}\selectfont
\renewcommand{\arraystretch}{1.15}
\setlength{\tabcolsep}{6pt}
\begin{longtable}{|>{\raggedright\arraybackslash}p{0.15\textwidth}|>{\raggedright\arraybackslash}p{0.11\textwidth}|>{\raggedright\arraybackslash}p{0.10\textwidth}|>{\raggedright\arraybackslash}p{0.56\textwidth}|}
\hline
\textbf{Extending Use Case} & \textbf{Base Use Case} & \textbf{Condition} & \textbf{Added Behavior} \\ \hline
\endhead
None & None & None & The proposals do not name a distinct extending use case for US-39. \\ \hline
\end{longtable}
\endgroup
\par\medskip

\hypertarget{sec:a4-model-notes}{}
\addcontentsline{toc}{subsubsection}{A4. Model Notes}
\subsubsection*{A4. Model Notes}
\par\medskip

\begin{enumerate}
\item UC-33 is derived solely from US-39; no I-xx, E-xx, or G-xx identifier is introduced.
\item The User actor is taken directly from the US-39 story wording "As a user".
\item No supporting system actor is drawn because the proposals name no external system for communication settings.
\item Settings is an aggregate hub in \texttt{settings.md}; Communication Settings is treated as a section of that hub, not as a separate use case.
\end{enumerate}
\par\medskip

\hypertarget{sec:part-b-use-case-specification}{}
\addcontentsline{toc}{subsection}{Part B --- Use Case Specification}
\subsection*{Part B --- Use Case Specification}
\par\medskip

\hypertarget{sec:b1-brief-description}{}
\addcontentsline{toc}{subsubsection}{B1. Brief Description}
\subsubsection*{B1. Brief Description}
\par\medskip

The user controls who may communicate with them and which interactions are permitted. US-39 requires block and mute controls for messaging, friend requests, and room invitations, restriction of direct messages to confirmed friends, and the ability to change communication and privacy settings at any time. The successful outcome is a saved preference set that the backend enforces so that disallowed interactions are prevented.
\par\medskip

\texttt{settings.md} groups these requirements into the Communication section of Settings, with dropdowns for direct messages, friend requests, and room invitations plus lists for muted and blocked members. \texttt{friend.md}, \texttt{friend\_profile.md}, and \texttt{group\_detail.md} describe the surfaces where the restrictions become visible, such as messaging, friend requests, and group member menus.
\par\medskip

\hypertarget{sec:b2-preconditions}{}
\addcontentsline{toc}{subsubsection}{B2. Preconditions}
\subsubsection*{B2. Preconditions}
\par\medskip

\begin{enumerate}
\item The user is authenticated with an active session (\texttt{NFR-SEC-1}).
\item The user can reach Settings from the main navigation; \texttt{screen\_flow.md} defines \texttt{Home -$>$ Settings}.
\item The Settings area and its Communication section are available to the user.
\item Existing communication preferences and any muted or blocked member lists can be read by the system.
\end{enumerate}
\par\medskip

\hypertarget{sec:b3-postconditions}{}
\addcontentsline{toc}{subsubsection}{B3. Postconditions}
\subsubsection*{B3. Postconditions}
\par\medskip

\begin{enumerate}
\item \textbf{Success:} The user's communication settings are stored with the selected direct-message, friend-request, and room-invitation values.
\item \textbf{Success:} The muted and blocked member lists reflect the user's final Mute, Unmute, Block, and Unblock actions.
\item \textbf{Success:} The backend prevents interactions the user does not permit (\texttt{FR-39.5}).
\item \textbf{Failure:} No setting change is persisted, and the user receives an appropriate error message.
\end{enumerate}
\par\medskip

\hypertarget{sec:b4-main-success-scenario-basic-flow}{}
\addcontentsline{toc}{subsubsection}{B4. Main Success Scenario (Basic Flow)}
\subsubsection*{B4. Main Success Scenario (Basic Flow)}
\par\medskip

\begin{enumerate}
\item The user opens \textbf{Settings} from the main navigation.
\item The system opens the Settings hub and highlights the \textbf{Communication} section.
\item The system renders \texttt{CommunicationSettingsForm} with the current \texttt{CommunicationPreference} values.
\item The user selects a value in \texttt{DirectMessageRestrictionSelect}: \textbf{Everyone} or \textbf{Confirmed friends only}.
\item The user selects a value in \texttt{FriendRequestPermissionSelect}: \textbf{Allow} or \textbf{Reject incoming friend requests}.
\item The user selects a value in \texttt{RoomInvitationPermissionSelect}: \textbf{Allow} or \textbf{Reject room invitations}.
\item The user adds or removes members through \texttt{MutedUsersList} and \texttt{BlockedUsersList}, using \textbf{Mute}, \textbf{Unmute}, \textbf{Block}, and \textbf{Unblock} as applicable.
\item The user selects \textbf{Save Communication}.
\item The frontend submits the selected values and the backend stores \texttt{directMessageRestriction}, \texttt{friendRequestsEnabled}, \texttt{roomInvitationsEnabled}, \texttt{mutedUserIds}, and \texttt{blockedUserIds}.
\item The backend applies the saved restrictions to relevant reads and interactions.
\item The system displays a success confirmation and returns the Communication section to a saved state.
\end{enumerate}
\par\medskip

\hypertarget{sec:b5-alternative-flows}{}
\addcontentsline{toc}{subsubsection}{B5. Alternative Flows}
\subsubsection*{B5. Alternative Flows}
\par\medskip

\hypertarget{sec:a1-update-privacy-settings}{}
\paragraph*{A1. Update privacy settings}
\par\medskip

The user opens the \textbf{Privacy} section instead of \textbf{Communication} and changes visibility, discoverability, or online-status controls under \texttt{FR-39.2}; \texttt{AC-39.3} permits communication and privacy settings to be changed at any time.
\par\medskip

\hypertarget{sec:a2-revoke-an-existing-restriction}{}
\paragraph*{A2. Revoke an existing restriction}
\par\medskip

At basic flow step 7, the user selects \textbf{Unmute} or \textbf{Unblock}. The corresponding member is removed from \texttt{MutedUsersList} or \texttt{BlockedUsersList} before the user saves.
\par\medskip

\hypertarget{sec:a3-open-the-reporting-flow}{}
\paragraph*{A3. Open the reporting flow}
\par\medskip

The user selects \texttt{ReportBehaviorLink}, labelled \textbf{Report inappropriate behavior}, to start reporting from Settings. This use case ends when the reporting flow opens; no moderator outcome is promised.
\par\medskip

\hypertarget{sec:b6-exception-flows}{}
\addcontentsline{toc}{subsubsection}{B6. Exception Flows}
\subsubsection*{B6. Exception Flows}
\par\medskip

\hypertarget{sec:e1-unauthenticated-access}{}
\paragraph*{E1. Unauthenticated access}
\par\medskip

The system redirects the user to Login before showing protected Settings content (\texttt{NFR-SEC-1}).
\par\medskip

\hypertarget{sec:e2-unsupported-selection}{}
\paragraph*{E2. Unsupported selection}
\par\medskip

The system keeps the user's input, displays an inline validation error through \texttt{SettingsErrorAlert}, and does not persist an unsupported value.
\par\medskip

\hypertarget{sec:e3-save-failure}{}
\paragraph*{E3. Save failure}
\par\medskip

The backend rejects the save. The system preserves the edited values, explains the failure, and allows the user to retry.
\par\medskip

\hypertarget{sec:e4-disallowed-interaction-after-save}{}
\paragraph*{E4. Disallowed interaction after save}
\par\medskip

A blocked or restricted member attempts messaging, a friend request, or a room invitation. The backend denies the interaction under \texttt{FR-39.5}; the UI may explain the restriction with \texttt{BlockedUserState} or \texttt{ConversationRestrictionNotice}.
\par\medskip

\hypertarget{sec:b7-business-rules}{}
\addcontentsline{toc}{subsubsection}{B7. Business Rules}
\subsubsection*{B7. Business Rules}
\par\medskip

\begin{enumerate}
\item The system shall allow a user to control communication settings (\texttt{FR-39.1}) and privacy settings (\texttt{FR-39.2}).
\item Users can mute or permanently block specific members from messaging, friend requests, or room invitations (\texttt{AC-39.1}, \texttt{FR-39.3}).
\item Users can restrict direct messages to confirmed friends (\texttt{AC-39.2}, \texttt{FR-39.4}).
\item \texttt{DirectMessageRestrictionSelect} = \textbf{Everyone / Confirmed friends only}, \texttt{FriendRequestPermissionSelect} = \textbf{Allow / Reject incoming friend requests}, \texttt{RoomInvitationPermissionSelect} = \textbf{Allow / Reject room invitations}, and mute/block lists = \texttt{MutedUsersList} / \texttt{BlockedUsersList} in \texttt{settings.md}.
\item Communication and privacy settings shall prevent interactions the user does not permit, with backend enforcement (\texttt{FR-39.5}, \texttt{AC-39.3}).
\item Users can change communication and privacy settings at any time (\texttt{AC-39.3}).
\item Friend requests, messages, and room invitations must respect the user's privacy choices (\texttt{NFR-PRIV-2}, \texttt{NFR-PRIV-6}, \texttt{NFR-PRIV-7}).
\end{enumerate}
\par\medskip

\hypertarget{sec:b8-special-requirements}{}
\addcontentsline{toc}{subsubsection}{B8. Special Requirements}
\subsubsection*{B8. Special Requirements}
\par\medskip

\begingroup
\small
\renewcommand{\baselinestretch}{1.0}\selectfont
\renewcommand{\arraystretch}{1.15}
\setlength{\tabcolsep}{6pt}
\begin{longtable}{|>{\raggedright\arraybackslash}p{0.14\textwidth}|>{\raggedright\arraybackslash}p{0.79\textwidth}|}
\hline
\textbf{ID} & \textbf{Requirement} \\ \hline
\endhead
NFR-SEC-1 & Authentication is required before protected Settings resources are accessible. \\ \hline
NFR-SEC-3 & Sensitive communication is encrypted using TLS 1.2 or higher. \\ \hline
NFR-SEC-6 & Unauthorized access is prevented through server-side authorization checks. \\ \hline
NFR-SEC-8 & The interface and backend protect against SQL injection, XSS, and CSRF. \\ \hline
NFR-PRIV-2 & Users can allow or reject incoming friend requests. \\ \hline
NFR-PRIV-6 & Users can control whether other users can send messages or room invitations. \\ \hline
NFR-PRIV-7 & Privacy settings apply consistently to forum posts, comments, member lists, and search results. \\ \hline
Design constraint & Restrictions must be enforced on the backend, not only by hiding UI controls (\texttt{FR-39.5}). \\ \hline
\end{longtable}
\endgroup
\par\medskip

\hypertarget{sec:b9-extension-points}{}
\addcontentsline{toc}{subsubsection}{B9. Extension Points}
\subsubsection*{B9. Extension Points}
\par\medskip

\begingroup
\small
\renewcommand{\baselinestretch}{1.0}\selectfont
\renewcommand{\arraystretch}{1.15}
\setlength{\tabcolsep}{6pt}
\begin{longtable}{|>{\raggedright\arraybackslash}p{0.23\textwidth}|>{\raggedright\arraybackslash}p{0.70\textwidth}|}
\hline
\textbf{Extension Point} & \textbf{Location in Flow of Events} \\ \hline
\endhead
Privacy section & Basic flow step 2; alternative A1 \\ \hline
Restriction revocation & Basic flow step 7; alternative A2 \\ \hline
Report entry & Communication section through \texttt{ReportBehaviorLink}; alternative A3 \\ \hline
Backend enforcement & Basic flow steps 9-10; exception E4 \\ \hline
Save failure recovery & Basic flow step 8; exception E3 \\ \hline
\end{longtable}
\endgroup
\par\medskip

\hypertarget{sec:b10-interface-and-data-notes}{}
\addcontentsline{toc}{subsubsection}{B10. Interface and Data Notes}
\subsubsection*{B10. Interface and Data Notes}
\par\medskip

\begin{itemize}
\item \textbf{Screen and feature flow.} \texttt{screen\_flow.md} defines \texttt{Home -$>$ Settings} and \texttt{Settings -$>$ Profile} / \texttt{Settings -$>$ Preferences}. \texttt{settings.md} recommends a Settings hub with section navigation and the route \texttt{/settings/communication}; the route is a recommendation, not an endpoint.
\item \textbf{Confirmed UI components.} \texttt{settings.md} names \texttt{CommunicationSettingsForm}, \texttt{DirectMessageRestrictionSelect}, \texttt{FriendRequestPermissionSelect}, \texttt{RoomInvitationPermissionSelect}, \texttt{MutedUsersList}, \texttt{BlockedUsersList}, \texttt{MuteUserButton}, \texttt{UnmuteUserButton}, \texttt{BlockUserButton}, \texttt{UnblockUserButton}, \texttt{ReportBehaviorLink}, and \texttt{SaveCommunicationSettingsButton}.
\item \textbf{Exact settings values.} \texttt{DirectMessageRestrictionSelect} = \textbf{Everyone / Confirmed friends only}; \texttt{FriendRequestPermissionSelect} = \textbf{Allow / Reject incoming friend requests}; \texttt{RoomInvitationPermissionSelect} = \textbf{Allow / Reject room invitations}. \texttt{MutedUsersList} shows muted members and mute scope; \texttt{BlockedUsersList} shows permanently blocked users.
\item \textbf{Exact actions and labels.} \texttt{Mute}, \texttt{Unmute}, \texttt{Block}, \texttt{Unblock}, \textbf{Save Communication}, and \textbf{Report inappropriate behavior} come from \texttt{settings.md}. \textbf{Save Changes}, \textbf{Discard}, and \textbf{You have unsaved changes} are recommended shell controls in \texttt{settings.md}.
\item \textbf{Data fields.} \texttt{CommunicationPreference} contains \texttt{directMessageRestriction}, \texttt{friendRequestsEnabled}, \texttt{roomInvitationsEnabled}, \texttt{mutedUserIds}, and \texttt{blockedUserIds} (\texttt{settings.md} information model).
\item \textbf{Enforcement surfaces.} \texttt{friend.md} defines \texttt{ConversationRestrictionNotice} with "Direct messages are limited to confirmed friends."; \texttt{friend\_profile.md} defines \texttt{BlockedUserState} with "You cannot interact with this member."; \texttt{group\_detail.md} defines \texttt{GroupAccessDeniedPanel} for private groups. These are UI feedback, while \texttt{FR-39.5} requires backend prevention.
\item \textbf{API and recommendation boundary.} Not specified in proposals for communication settings; \texttt{API.md} defines no communication-settings endpoint, payload, status code, or error code. Do not assign forum codes such as \texttt{48010 FORBIDDEN} or database codes such as \texttt{45000 UNKNOWN ERROR} to settings actions. Routes, component names, \texttt{SettingsSuccessToast}, and \texttt{SettingsErrorAlert} are UI/UX recommendations in \texttt{settings.md}; the requirements are \texttt{AC-39.1}-\texttt{AC-39.3} and \texttt{FR-39.1}-\texttt{FR-39.5}. \texttt{NFR-PRIV-2}, \texttt{NFR-PRIV-6}, and \texttt{NFR-PRIV-7} require consistent control of friend requests, messages, room invitations, and privacy across posts, comments, member lists, and search results.
\end{itemize}
\par\medskip

\hypertarget{sec:b11-source-alignment-notes}{}
\addcontentsline{toc}{subsubsection}{B11. Source Alignment Notes}
\subsubsection*{B11. Source Alignment Notes}
\par\medskip

\begin{itemize}
\item \texttt{userstory.md} sections 11.2-11.3 and \texttt{functional\_requirement.md} US-39 agree on the three AC goals and five FR requirements; \texttt{settings.md} supplies the concrete dropdown values and mute/block lists.
\item No communication-settings or behavioral-report API endpoints are defined for this use case: \texttt{API.md} has no communication-settings endpoint, and its only report endpoint is the forum-specific \texttt{POST /forum/reports}, which does not cover profiles, messages, or users.
\item The \texttt{ReportBehaviorLink} labelled \textbf{Report inappropriate behavior} is an entry point only; \texttt{settings.md} is a UI/UX recommendation, and report UI must not promise a moderator outcome. Submission feedback is covered by UC-34 and \texttt{ReportSubmittedNotice}.
\item \texttt{FR-39.5} requires backend enforcement, but no backend contract or error-code contract is defined for communication restrictions; \texttt{friend.md}, \texttt{friend\_profile.md}, and \texttt{group\_detail.md} have no friend/profile/group API contract and contribute only interaction-feedback surfaces.
\item No distinct include, extend, or specialization use case is named for US-39, so no I-xx, E-xx, or G-xx identifier is created.
\item Forum-specific content reporting is covered by PA4 \texttt{UC-20-report-moderate-forum-content.md}; this use case covers communication controls under US-39 and points behavioral reports to UC-34 without duplicating moderation decisions.
\end{itemize}
\par\medskip

\par\medskip
\hypertarget{sec:ui-prototype-specimens}{}
\addcontentsline{toc}{subsection}{UI Prototype Specimens}
\subsection*{UI Prototype Specimens}
\par\medskip
\begin{figure}[H]
\centering
\includegraphics[width=0.75\textwidth,height=0.42\textheight,keepaspectratio]{"../../prototype/settings_04a_settings_notifications_selected.png"}
\par\vspace{2pt}
{\footnotesize \textbf{UI Prototype:} Settings: 04a Settings - Notifications selected}
\end{figure}
\par\smallskip

\end{document}
